%% title: Long-tail rebalancing inverts under shift; cRT removes the trade
%% target: Sections/Results.tex
%% anchor: A representation trained on 43 times more images of the same domain is better before adaptation, and the adaptation recipe closes the difference.
%% action: after
%% why: Every push of probability mass toward rare species lowers accuracy under shift monotonically, invisibly on the test fold; the cost grows with capacity; rebalancing only the classifier keeps most of the in-distribution gain without it. Same shape for the per-species image cap.
%% source: paper/DRAFT.md §4.13, §4.13.1; journal/2026-08-01-imbalance-methods-bench.md; journal/2026-10-22-rescoring-in-the-mini-metrics-convention.md (in-distribution calibrated; full-trap rows); data/rescore/trap-grouped.csv (L0 0.3490, base 0.2836, L1 0.2241, L2 0.1834); probe-grouped.csv (base-probef 0.2875, L4 0.3433, F1 0.4909, L5 0.5364, cap250 0.2734, cap500 0.3202, cap1000 0.3597, c1000b 0.3436); probe-ho-grouped.csv (base 0.2394, F1 0.3900, L5 0.4384, cap 0.2285 / 0.2613 / 0.2745 / 0.2863)
%% check: Paste after step 18. | Decision sheet row 3 recommends Asger's no-oversampling plus cRT as a reported stage; this piece is the evidence. | L0 was not scored on the probe set and L4 not on the full trap set, so cRT is compared with the oversampled baseline on the probe set only. | Uncapped in-distribution: 0.9135 (this baseline) or 0.9148 (DRAFT §3.1 ladder); use one. | Balanced softmax equals logit adjustment at tau 1 and EMLA is a logit adjustment: say how EMLA relates to L1 once Asger's runs are scored on traps (decision sheet row 4).
%% status: ready
\subsection{Long-tail rebalancing inverts under shift}\label{sec:results_longtail}
Long-tailed recognition is usually validated on held-out splits of the training distribution.
\Cref{tab:longtail} crosses square-root oversampling with balanced softmax, which is logit adjustment at $\tau = 1$ \textbf{(REF: Ren et al. 2020; Menon et al. 2021)}, on one backbone with recipe L.
Ordered by how hard they push probability mass toward rare species, the four cells lose accuracy on the trap images at every step, 16.6 points from first to last against a 1.0-point difference between identical trainings.
On the test fold they do not order at all.
What a model learns from a species' few dozen photographs is largely about those photographs (one photographer, background and camera), so tail reweighting up-weights the least transferable part of the training signal, which the test fold, sharing those artefacts, cannot see.

\begin{table}[htb]
    \centering
    \caption{Rebalancing the long tail (EfficientNetV2-S, recipe L except for the factor varied). Test fold: zero threshold and calibrated. Traps: all 47,905 crops; probe: the probe set; both calibrated on 10\,\% of the (trap, night) groups.}
    \label{tab:longtail}
    \begin{tabular}{lcccc}
        \toprule
        & \multicolumn{2}{c}{test fold} & traps & probe \\
        rebalancing & zero & calibrated & & \\
        \midrule
        none (L0) & 0.8949 & 0.8828 & 0.349 & -- \\
        square-root oversampling & 0.9135 & 0.9117 & 0.284 & 0.288 \\
        balanced softmax (L1) & 0.8970 & 0.9028 & 0.224 & -- \\
        both (L2) & 0.8689 & 0.8877 & 0.183 & -- \\
        classifier only, cRT (L4) & 0.9068 & 0.9045 & -- & 0.343 \\
        \bottomrule
    \end{tabular}
\end{table}

The cost grows with capacity.
At 198\,M parameters, removing oversampling costs 1.6 points on the test fold and gains 4.6 on the probe set and 4.8 on its held-out species.
Constraints on the objective weaken with capacity (at 198\,M, trap-like augmentation no longer costs anything on the test fold, and marginal supervision no longer helps under shift); oversampling changes which data the model sees, and a larger model fits the reweighted distribution, and its bias, better.

The trade is not intrinsic.
Training the representation on the natural distribution and rebalancing only the classifier on the frozen backbone (cRT, \cref{sec:selftraining}) recovers 64\,\% of oversampling's gain on the test fold (75\,\% calibrated) and scores 5.6 points above the oversampled model on the probe set.
The damage is done in the representation.

The number of images per species behaves the same way.
Capping training at 250, 500 and 1,000 images per species instead of about 2,000 raises the test fold monotonically with the cap (0.878, 0.896, 0.906, 0.914), while the probe set and its held-out species peak at 1,000 (probe 0.273, 0.320, 0.352 and 0.288 for the full cap; mean of two runs at 1,000).
Images of species that already have a thousand buy accuracy on the saturated axis and cost it under shift.
